% --------------------------------------------------------------------------
\section{Augmented Clustering with Gradient Fingerprints}
\label{sec:ch4_gradfp}

Representation-space clustering alone may not fully capture adversarial difficulty.
Two adversarial examples can be close in feature space while inducing different
optimization behavior. To capture this additional signal, AttackDRO++ constructs
an augmented clustering feature that combines representation similarity, class
information, and gradient-level training dynamics.

For an adversarial example \(\tilde{\mathbf{x}}_i=\mathbf{x}_i^{\mathrm{adv}}\), let
\[
z_i = f_\theta(\tilde{\mathbf{x}}_i),
\qquad
h_i = \phi_\theta(\tilde{\mathbf{x}}_i).
\]
The logits \(z_i=f_\theta(\tilde{\mathbf{x}}_i)\) are used to compute prediction
probabilities and gradient-fingerprint terms, while
\(h_i=\phi_\theta(\tilde{\mathbf{x}}_i)\) provides the representation component of
the clustering feature. Let
\[
p_i = \mathrm{softmax}(z_i)
\]
be the predicted class distribution.

The first component is the normalized representation feature,
\[
\bar{h}_i = \operatorname{norm}(h_i),
\]
which captures semantic similarity in the model feature space.

The second component is a class-aware scalar label feature. When label information
is enabled, the implementation appends the normalized scalar class-index feature
\(\tilde{y}_i \in [0,1]\) with weight \(\lambda_{\mathrm{lbl}}\):
\[
\lambda_{\mathrm{lbl}}\tilde{y}_i,
\]
which provides weak label structure during clustering without using a one-hot label
vector in the clustering feature.

The third component is the gradient fingerprint. For a linear classifier head
\[
f_\theta(\cdot) = W\phi_\theta(\cdot) + b,
\qquad
W \in \mathbb{R}^{C \times d_h},
\]
the per-example gradient of the cross-entropy loss with respect to the classifier
weight matrix is
\[
\nabla_W \ell(f_\theta(\tilde{\mathbf{x}}_i),y_i)
=
\bigl(p_i-\mathrm{onehot}(y_i)\bigr)\otimes h_i
\in \mathbb{R}^{C\times d_h}.
\]
We use this quantity as a gradient proxy and denote it by
\[
G_i =
\bigl(p_i-\mathrm{onehot}(y_i)\bigr)\otimes h_i.
\]
The one-hot vector appears here only in the cross-entropy gradient proxy; it is
not appended as the label component of \(\psi_i\).

Intuitively, \(G_i\) describes the training direction induced by sample \(i\).
Thus, two adversarial examples may be grouped together not only because they are
close in representation space, but also because they induce similar optimization
updates.

The raw gradient proxy has dimension \(C d_h\). For CIFAR-10 with a ResNet-18
backbone, this corresponds to \(10 \times 512 = 5120\) dimensions. To make the
fingerprint more compact and efficient for clustering, we use a fixed random
projection matrix
\[
P \in \mathbb{R}^{C d_h \times d_{\mathrm{proj}}},
\]
and define the projected gradient fingerprint as
\[
g_i = \mathrm{vec}(G_i)P
\in \mathbb{R}^{d_{\mathrm{proj}}}.
\]
In our default setting, \(d_{\mathrm{proj}}=128\). The projection is fixed throughout
training and introduces no additional learnable parameters.

The final augmented clustering feature is
\[
\psi_i =
\left[
\operatorname{norm}(h_i)
\;\Vert\;
\lambda_{\mathrm{lbl}}\tilde{y}_i
\;\Vert\;
\lambda_{\mathrm{grad}}\operatorname{norm}(g_i)
\right].
\]

The three components play complementary roles. Representation features describe where
an example lies in feature space, whereas gradient fingerprints describe how the
classifier head would need to change to classify that example correctly. Combining
representation, scalar label, and gradient information makes the clustering signal
both representation-aware and optimization-aware. This allows AttackDRO++ to discover
latent adversarial groups that reflect both structural similarity and optimization
difficulty.
